% บทคัดย่อภาษาไทย
% หมายเหตุ: ไม่ควรเกิน 300 คำ
% ควรครอบคลุม: วัตถุประสงค์ → วิธีการ → ผลลัพธ์ → การประยุกต์ใช้
%
% *** ลบข้อความตัวอย่างด้านล่างทั้งหมดแล้วเขียนบทคัดย่อของตนเอง ***

โครงงานนี้มีวัตถุประสงค์เพื่อพัฒนาเว็บแอปพลิเคชันต้นแบบสำหรับระบบแนะนำหนังสือการ์ตูนและออกแบบโมเดลระบบแนะนำที่วิเคราะห์พฤติกรรมผู้ใช้หลายรูปแบบเพื่อเพิ่มความแม่นยำในการแนะนำสำหรับการเช่าอ่านโดยประยุกต์ใช้โมเดล MB-CGCN (Multi-Behavior Collaborative Graph Convolutional Network) วิเคราะห์พฤติกรรมผู้ใช้ในหลายมิติได้แก่การคลิกอ่านเนื้อหาการเพิ่มเข้าตะกร้าและการเช่าอ่าน (พฤติกรรมเป้าหมาย) พร้อมทั้งพัฒนาเว็บแอปพลิเคชันต้นแบบเพื่อแสดงผลและทดสอบประสิทธิภาพของโมเดลและแก้ไขปัญหา Cold-Start ด้วยการให้ผู้ใช้ใหม่เลือกหนังสือการ์ตูนที่สนใจระหว่างการสมัครสมาชิกรวมถึงการอัปเดตผลการแนะนำแบบทันทีเมื่อผู้ใช้เพิ่มหนังสือลงในตะกร้าหรือคลิกอ่านเนื้อหา

โดยทำการประเมินและเปรียบเทียบประสิทธิภาพกับระบบแนะนำแบบพฤติกรรมเดียวด้วยตัวชี้วัด Ranking Metrics ได้แก่ Recall@K และ NDCG@K เพื่อวิเคราะห์ว่าการใช้ข้อมูลพฤติกรรมหลายรูปแบบช่วยลดปัญหา data sparsity และสะท้อนความสนใจของผู้ใช้ได้ดีขึ้นหรือไม่ซึ่งสามารถนำไปประยุกต์ใช้ในธุรกิจเช่าหนังสือการ์ตูนดิจิทัลเพื่อเพิ่มประสิทธิภาพการแนะนำเนื้อหาในเชิงพาณิชย์

